\scanpage{69}
Dans le cas de \deleted{\unsure{Quat($S$)}} \deleted{Les foncteurs en sens inverse s'obtiennent}
\[
\operatorname{Quat}(\unsure{\mathcal O_S|S|\mathscr S})\longrightarrow\operatorname{For}(\mathrm{SL}(2)_S)
\]
on trouve
\[
M\longrightarrow\widetilde G\simeq\deleted{\unsure{\check{\mathbb V}|\check{\mathrm V}}(M)^*/\unsure{G_m}}\operatorname{Ker}\bigl(\check{\mathbb V}(M)^{\unsure{*|1}}\xrightarrow{\det}G_m\bigr)\tag*{(29)}
\]
\inserted{$=$ \uline{norme réduite}}

Pour les foncteurs de $\operatorname{For}(\mathrm{SL}(2)_S)$ vers les autres catégories, on peut soit passer par $\operatorname{Tor}(\mathrm{GP}(1))$ ou par $\operatorname{For}(\mathrm{GP}(1)_S)$, soit dans certains cas procéder directement, p.ex.
\[
\begin{aligned}
\operatorname{For}(\mathrm{SL}(2)_S)&\longrightarrow\operatorname{For}(\mathfrak{sl}(2)_S)\\
\widetilde G&\longrightarrow\widetilde{\mathfrak g}=\underline{\operatorname{Lie}}(\widetilde G)
\end{aligned}\tag*{(30)}
\]
\[
\widetilde G\longrightarrow X\simeq\underline{\operatorname{Bor}}(\widetilde G)\tag*{(31)}
\]

Pour la relation \deleted{\unsure{triple} entre $G,\widetilde G,\unsure{(E,q,\omega)|M,q},\mathfrak g,\widetilde{\mathfrak g}$} \inserted{\unsure{des ob}} \deleted{\unsure{on}} $G,\widetilde G,\mathfrak g,\widetilde{\mathfrak g},(E,q,\omega)$, c'est $M$ qui \unsure{semble|donne} peut être la meilleure clef pour les obtenir tous -- ainsi que $X$ et $(P,C)$, tant qu'à faire ! -- et pour comprendre les relations entre tous ces objets. On verra ainsi que $\mathfrak g$ et $\widetilde{\mathfrak g}$ sont des \unsure{modules duels|modules duaux} l'une de l'autre, que $E\simeq\widetilde{\mathfrak g}$ canoniquement i.e. qu'en \unsure{parti-|parti.}%
% archive sheet number: 68
